\documentclass[10pt,a4paper]{amsart}
\usepackage[margin=19mm]{geometry}
\usepackage{amsmath,amssymb,mathtools}
\usepackage[scr=rsfs,cal=euler]{mathalfa}
\usepackage{xfrac}
\usepackage[unicode,hidelinks,bookmarks=false]{hyperref}
\newcommand{\ie}{i.e.\ }
\newcommand{\cf}{cf.\ }
\newcommand{\resp}{respectively\ }
\newcommand{\Subsection}{\subsection}
\providecommand{\nobreakdash}{\nobreak\hbox{-}}
\setlength{\emergencystretch}{2em}
\multlinegap0pt
\usepackage{bm}
\usepackage{bbm}
\usepackage{dsfont}
\usepackage{xspace}
\usepackage{cancel}
\usepackage{mathtools}
\usepackage{amsthm}
\makeatletter
\@ifundefined{thm}{\newtheorem{thm}{Thm}}{}
\@ifundefined{dfn}{\newtheorem{dfn}{Definition}}{}
\@ifundefined{lem}{\newtheorem{lem}[thm]{Lemma}}{}
\@ifundefined{prop}{\newtheorem{prop}[thm]{Proposition}}{}
\makeatother
\newcommand{\R}{\mathbb{R}}
\newcommand{\del}{\partial}
\title[Bi-contact structures with symmetry: local normal forms]{Bi-contact structures with symmetry: local normal forms}
\author{Connor Jackman}
\date{Source: arXiv:2508.00802v2 (CC BY 4.0)}
\begin{document}
\maketitle

\noindent\emph{Source and licence.} Bi-contact structures with symmetry: local normal forms. Version arXiv:2508.00802v2 (CC BY 4.0), distributed under the Creative Commons Attribution 4.0 International licence, as recorded in the arXiv licence metadata for this exact version and shown on the abstract page https://arxiv.org/abs/2508.00802v2. Full rights notice: licenses/arxiv-2508.00802-CC-BY-4.0.txt.\par\smallskip

\noindent\textit{Editorial scope.} Counted unit: the Prop whose source label is \texttt{prop:SymmsDepGen} in the article (source0.tex lines 745--753, source label \texttt{prop:SymmsDepGen}), delivered together with the article's own complete original proof of it (source0.tex lines 754--841, one \texttt{proof} environment of 88 lines). No mathematics is written, repaired or re-derived: statement and proof are the article's own text line for line. Required context retained uncounted: def:BalGenInv (lines 366--377); eq:GenBalStrEq (lines 522--528); eq:Ric (lines 532--534); prop:DepGenlnf (lines 731--734). Every definition, display and label of those lines is retained unchanged. Omitted from this package and not redistributed: 1--365, 378--521, 529--531, 535--730, 735--744, 842--1002. Nothing inside a retained excerpt is omitted or altered: every retained body is delivered exactly as the article's lines are written, so no omission receipt of any kind accompanies this package. Citations: the retained excerpts carry no citation command at all, so no bibliography entry of the article is transcribed and no reference is redistributed. Reference adaptation: none was needed and none was made; every reference inside the delivered unit resolves inside it, so no unresolved reference or citation is produced. Printed slips kept literal: no printed word, symbol, hypothesis or number of the counted statement or of its proof was altered. Layout and class: the article's own class is \texttt{amsart} with packages including babel, amsmath, amssymb, latexsym, amsthm, enumerate, mathrsfs, ulem, epsfig, graphicx, wrapfig, color, tikz, hyperref; the wrapper is the lane's pinned \texttt{amsart} editorial wrapper at 10pt on A4, which restates the theorem-like environments the retained text needs and adds the article's own macro definitions that the retained text needs (\texttt{\textbackslash R}, \texttt{\textbackslash del}), with extra wrapper packages (bm, bbm, dsfont, xspace, cancel, mathtools). The delivered numbering is the unit's own. Assets: no retained span contains a figure environment, a graphics-inclusion command, a PDF-page inclusion, a table or a TikZ picture; no image file is copied into this package, the package declares no asset dependency and no image file is redistributed with it. Licence: version arXiv:2508.00802v2 of this article is distributed under the Creative Commons Attribution 4.0 International licence, as recorded in the arXiv licence metadata for this exact version and shown on the abstract page https://arxiv.org/abs/2508.00802v2. A full rights notice is delivered at licenses/arxiv-2508.00802-CC-BY-4.0.txt. Characters: every retained excerpt is pure ASCII, so the delivered engine, XeLaTeX with its default Latin Modern text font, reports no missing character.
\begin{dfn}\label{def:BalGenInv}
{\em We call a pair $(\alpha, \tilde\alpha)$ of contact forms $\ker \alpha = \xi, \ker\tilde\alpha = \tilde \xi$ for $(M, \xi, \tilde \xi)_\pm$ {\em balanced} when:
    \[ \alpha\wedge d\alpha = \pm \tilde\alpha\wedge d\tilde\alpha. \]
    Supposing the axis has been oriented by $A = \text{span}_+\{ X\}$, we take the {\em normalized vector field} $X:M\to A$ spanning the axis through:
    \[ i_X(\alpha\wedge d\alpha) = \alpha\wedge\tilde\alpha \]
    for $\alpha, \tilde\alpha$ a choice of balanced contact forms. The corresponding one form along the axis we denote as, $ds : M\to A^*$,
    with $ds(X) = 1$. For a function $f:M\to \R$, we denote its {\em derivative along the axis} by:
    \[ f' := Xf :M\to \R.\]
    As well, we have an induced function, which we call the  {\em generating invariant}:
    \[ I:M\to \R, ~~~I := \frac{\alpha\wedge d\tilde\alpha + \tilde\alpha \wedge d\alpha}{\alpha\wedge d\alpha} \]
    where $\alpha, \tilde\alpha$ are balanced contact forms such that $i_X(\alpha\wedge d\alpha) = \alpha\wedge\tilde\alpha$.}
\end{dfn}

\begin{prop}
    Let $\alpha, \tilde\alpha$ be balanced contact forms for $(M, \xi, \tilde \xi)_\pm$ and $\beta$ a 1-form on $M$ with $\beta|_A = ds$. Then such contact forms satisfy:
    \begin{equation}\label{eq:GenBalStrEq}
        d\alpha = \tilde h \alpha\wedge\tilde\alpha + \tilde\alpha\wedge \beta + g \beta\wedge\alpha , ~~d\tilde\alpha = h \alpha\wedge\tilde\alpha + (I - g )\tilde\alpha\wedge \beta \pm  \beta\wedge\alpha
    \end{equation}
    for some functions $h, \tilde h,  g$ on $M$ (and where $I:M\to\R$ is the generating invariant from definition \ref{def:BalGenInv}). Note that $\alpha, \tilde\alpha$ are determined up to simultaneous scalings: $\alpha, \tilde\alpha\mapsto \lambda\alpha, \lambda\tilde\alpha$, and $\beta$ up to shifts $\beta\mapsto \beta + b_1\alpha + b_2\tilde\alpha$.
\end{prop}

    \begin{equation}\label{eq:Ric}
        \rho' = 1 + I\rho \pm \rho^2.
    \end{equation}

\begin{prop}\label{prop:DepGenlnf}
    For $(M, \xi, \tilde \xi)$ with $\mathcal{S}\ne 0$ and $I'(I)$ a function of $I$, there are local coordinates in which: $\xi = \{ dy = p e^{c(x,y)}dx \}, ~\tilde \xi = \{ dy = f(p) e^{c(x,y)}dx\}$,
    for some functions $c(x,y)$ and $f(p)$.
\end{prop}

\begin{prop}\label{prop:SymmsDepGen}
    If a pair $(M, \xi, \tilde \xi)$ with $\mathcal{S}\ne 0$ and $I'(I)\ne 0$ a function of $I$ admits an infinitesimal symmetry then either
    \begin{enumerate}
        \item there are local coordinates $(x,y,p)$ in which:
        \[ \xi = \{ dy = p dx\}, ~~\tilde \xi = \{ dy = f(p) dx\}, \]
        \item or, there are local coordinates $(x,y,p)$ in which:
    \[ \xi = \{  dy = \frac{p + a(y)}{p + b(y)} dx \} , ~~\tilde \xi = \{ dy = \frac{f(p) + a(y)}{f(p) + b(y)} dx \}. \]
    \end{enumerate}
\end{prop}

\begin{proof}
    We prolong in this case to consider the principal $\R^\times$ sub-bundle of the co-frame bundle on $M$
    \[ \R^\times \to B\to M \]
    consisting of co-frames $\alpha, \tilde\alpha, \beta = dI/I'$ on $M$, where the $\R^\times$ action is by simultaneous re-scalings $\alpha,\tilde\alpha \mapsto \lambda\alpha, \lambda\tilde\alpha$ on the balanced contact forms $\alpha, \tilde\alpha$. We first claim that we have on $B$ a canonical co-frame induced by the pair $(M,\xi, \tilde \xi)$ with $I'(I)\ne 0$. Let $\Theta, \tilde\Theta, \beta$ be the tautological 1-forms on $B$: $\Theta_b(v) = \alpha(\pi_*v), \tilde\Theta_b(v) = \tilde\alpha(\pi_*v)$, and $\beta = \pi^*\beta$ for $b = (\alpha, \tilde\alpha, \beta)\in B$ and $\pi:B\to M$. Then:
    \begin{lem}
         There is a unique connection 1-form $\nu$ on $B$ for which the co-frame $\Theta, \tilde\Theta, \beta, \nu$ satisfies the structure equations:
    \begin{equation}\label{eq:liftStrEq}
        d\Theta = \nu\wedge\Theta + \tilde\Theta\wedge \beta, ~~~d\tilde\Theta = \nu\wedge\Tilde\Theta + I\tilde\Theta\wedge\beta \pm \beta\wedge\Theta , ~~d\beta = 0, ~~d\nu = \Lambda \Theta\wedge\tilde\Theta,
    \end{equation}
    the coefficient $\Lambda:B\to \R$ being another invariant of $(M,\xi, \tilde \xi)$.
    \end{lem}
    \begin{proof}
        Indeed, for any choice of balanced contact forms with trivialization $B\cong M\times\R\ni (m, \lambda)$ through $\Theta = \lambda\alpha, \tilde\Theta = \lambda \tilde\alpha$, we take:
    \[ \nu = \frac{d\lambda}{\lambda} + h \alpha - \tilde h \tilde\alpha + g\beta,   \]
    where $\alpha, \tilde\alpha, \beta$ satisfy \eqref{eq:GenBalStrEq}.
    \end{proof}
    Now, having a canonical co-frame on the 4-dimensional $B$ with (at-least) one non-trivial invariant $I$, such a pair can admit at most a 3 dimensional symmetry algebra: only when $\Lambda(I)$ is dependent on $I$. It follows directly by differentiating \eqref{eq:liftStrEq} that:
    \[ d\Lambda = \Lambda_1 \Theta + \Lambda_2\tilde\Theta + I\Lambda\beta - 2\Lambda\nu \]
    so that $\Lambda$ is dependent on $I$ ( $I'(I)\beta = dI$ ) iff $\Lambda\equiv 0$. In case $\Lambda\equiv 0$, we may take a co-frame as in proposition~\ref{prop:DepGenlnf}, and directly compute that $\Lambda\equiv 0$ is equivalent to $\del_x\del_y c \equiv 0$, ie $c(x,y) = a(x) + b(y)$ is separable, and the normal form of proposition~\ref{prop:DepGenlnf} takes the form stated in item 1 above.

    
    Otherwise, when $\Lambda\ne 0$, we have two independent and non-trivial invariants, $I, \Lambda$, and (at most) a 2-dimensional symmetry algebra. As it turns out:
    \begin{lem}
        A pair $(M, \xi, \tilde\xi)$ with $I'(I)\ne 0$ and $\Lambda\ne 0$ admits at most a 1-dimensional symmetry algebra.
    \end{lem}
    \begin{proof}
        To realize a 2-dimensional symmetry algebra the additional invariants $\Lambda_1, \Lambda_2$ would need to be functionally dependent on $I, \Lambda$.
    One computes directly that:
    \[ d\Lambda_1 = \Lambda_{11}\Theta + \Lambda_{12}\tilde\Theta + (I\Lambda_1 \mp \Lambda_2) \beta - 3\Lambda_1\nu, ~~ d\Lambda_2 = \Lambda_{21}\Theta + \Lambda_{22} \tilde\Theta + (\Lambda_1 + 2I\Lambda_2)\beta - 3\Lambda_2\nu \]
where $\Lambda_{21} - \Lambda_{12} = 2 \Lambda^2$.
In particular, if $d\Lambda_1, d\Lambda_2 \in \langle dI, d\Lambda\rangle$ (and $\Lambda\ne 0$) then we find the contradiction:
\[ 2\Lambda^3 = \Lambda(\Lambda_{21} - \Lambda_{12}) = 3( \Lambda_1\Lambda_2 - \Lambda_1\Lambda_2) = 0. \]
    \end{proof}

     When $\Lambda\ne 0$, we may pass back down to $M$, where we have the co-frame:
\[ \theta = |\Lambda|^{1/2}\Theta, ~\tilde\theta = |\Lambda|^{1/2} \tilde\Theta, ~\beta = dI/I' \]
whose structure equations we compute as:
\begin{equation}\label{eq:DownEstr}
 d\theta = - \frac{n}{2}\theta\wedge\tilde\theta + \tilde\theta\wedge\beta + \frac{I}{2} \beta\wedge\theta, ~~d\tilde\theta = \frac{m}{2} \theta\wedge\tilde\theta + \frac{I}{2}\tilde\theta\wedge\beta \pm \beta \wedge \theta,~~
d\beta = 0
\end{equation}
for the invariants
\[ m = \Lambda_1/|\Lambda|^{3/2}, ~n = \Lambda_2/|\Lambda|^{3/2} \]
on $M$, having
\[ dm = m_1\theta + m_2\tilde\theta -  \left( \frac{Im}{2} \pm n\right)\beta, ~~ dn = n_1 \theta + n_2 \tilde\theta + \left( m + \frac{In}{2} \right)\beta,  \]
\[n_1 = m_2 + 2. \]
In general $I, m,n:M\to \R$ are independent, and such a pair with $\Lambda\ne 0$ admits no infinitesimal symmetries.  Such a pair admits an infinitesimal symmetry only when the differentials $\langle dm, dn, dm_1, dm_2, dI\rangle$ have rank two:
\begin{equation}\label{eq:Dep1SymCond}
    0 \equiv \begin{vmatrix}
    m_1 & m_2 \\ n_1 & n_2
\end{vmatrix},  ~~0 \equiv \begin{vmatrix}
    m_1 & m_2 \\ n_{21} & n_{22}
\end{vmatrix}, ~~0 \equiv \begin{vmatrix}
    m_1 & m_2 \\ m_{21} & m_{22}
\end{vmatrix}, ~~0 \equiv \begin{vmatrix}
    m_1 & m_2 \\ m_{11} & m_{12}
\end{vmatrix}.
\end{equation}
Supposing the condition \eqref{eq:Dep1SymCond} holds, we now proceed to determine the normal form of item 2, by seeking appropriate solutions of our Ricatti equation \eqref{eq:Ric}. Note that for any function $f:M\to \R$, writing $df = f_1\theta + f_2\tilde\theta + f'\beta$, then by differentiating and using \eqref{eq:DownEstr}, one obtains the relations:
\begin{equation}\label{eq:Dep1SymDerivs}
    (f_2)' - (f')_2 = f_1 + \frac{I}{2} f_2, ~~ (f')_1 - (f_1)'  =  \frac{I}{2}f_1 \pm f_2 , ~~ f_{21} - f_{12} = \frac{nf_1 - mf_2}{2}
\end{equation} 
for its higher derivatives.
From \eqref{eq:Dep1SymCond}, we set:
\[ \rho = m_2/m_1 = n_2/n_1, ~~\ell = n/m,  \]
and check from \eqref{eq:Dep1SymDerivs} that $\rho, \ell$ are solutions to \eqref{eq:Ric} and, for $d\rho = \rho_1\theta + \rho \rho_1\tilde\theta + \rho'\beta$, we have:
\[ \ell = \rho + 2\frac{\rho_1}{m}. \]
Note that we cannot have $\rho_1\equiv 0$ (ie $\ell\equiv\rho$) since differentiating $n = \rho m$, and using $n_1 = m_2 + 2$ one would then have:
\[ m_2 = \rho m_1 = n_1 = m_2 + 2, \]
which is impossible, so that $\rho_1\not\equiv 0$, and $\ell\not\equiv\rho$ are distinct solutions to \eqref{eq:Ric}. Then, using \eqref{eq:DownEstr} -- \eqref{eq:Dep1SymDerivs}, one computes:
\[ d\left( \frac{\theta + \rho\tilde\theta}{\rho_1} \right) = 0 , ~~~ d\left( m (\theta + \ell\tilde\theta) \right) = dy \wedge \left( m(\theta + \ell\tilde\theta) \right)  \]
so that for the integrable distributions $\theta + \rho\tilde\theta, \theta + \ell \tilde\theta$ we have integrating factors:
\[ \rho_1dy = \theta + \rho\tilde\theta, ~~ \frac{e^y}{m} dx = \theta + \ell\tilde\theta, \]
and the contact forms are given through:
\begin{equation}\label{eq:ContFormsDep}
    m(\ell - \rho) \theta = \ell m\rho_1 dy - \rho e^{y} dx, ~~m(\rho - \ell)\tilde\theta = m\rho_1dy - e^y dx.
\end{equation}
To give them the more explicit form of item 2 above, we note, from \eqref{eq:Dep1SymCond}, that all the invariants depend only on $y, I$, and for a function $f(y,I)$ we have:
\[ \del_yf = \rho_1f_1. \]
In particular, we compute that $(\rho_1\ell_1 m^2)' \equiv 0$, so that:
\[ m^2 \del_y\ell  = \rho_1\ell_1 m^2 = c(y) \]
is some function of $y$. Now, since $\rho, \ell$ are solutions to \eqref{eq:Ric}, they have the form:
\[ \frac{\ell - r_1}{\ell - r_2} = a(y) e^\lambda, ~~\frac{\rho - r_1}{\rho - r_2} = b(y) e^\lambda \]
where $r_j(I)$ are any pair of solutions to \eqref{eq:Ric} depending only on $I$ (and $\lambda(I) = \pm \int \frac{r_1 - r_2}{I'} dI$). Computing then the partials $\del_y\ell = c(y)/m^2, \del_y\rho = (\rho_1)^2$ and substituting into \eqref{eq:ContFormsDep} we arrive to the contact 1-forms:
\[ dY = \gamma(y) dy = \frac{r_1(I) - a(Y) e^{\lambda(I)}r_2(I) }{r_1(I) - b(Y) e^{\lambda(I)}r_2(I)} dx, ~~dY = \frac{1 - e^{\lambda(I)}a(Y)}{1 - e^{\lambda(I)}b(Y)}dx \]
for some function $\gamma(y)$,
which is in the stated form of item 2 above (set $P = \frac{r_1}{e^\lambda r_2}$ and $f(P) = e^{-\lambda}$).
\end{proof}

\end{document}
